\chapter{Foundations of language models}
\label{ch:foundations}

\begin{quote}
\emph{``A language model is a probability distribution over sequences of tokens. Everything else --- chatbots, code generation, summarization --- is a consequence of that single idea.''}
\end{quote}

% ============================================================
\section{What is a language model?}

A language model assigns a probability to a sequence of tokens $w_1, w_2, \ldots, w_n$:
\[
P(w_1, w_2, \ldots, w_n) = \prod_{i=1}^{n} P(w_i \mid w_1, \ldots, w_{i-1})
\]

This factorization says: predict each token given everything that came before it. A model that does this well can generate coherent text, translate languages, answer questions, and write code --- because all of these tasks reduce to ``what token comes next?''

\begin{definition}[Language model]
A function $P_\theta$ parameterized by $\theta$ that maps a sequence of tokens to a probability. Training adjusts $\theta$ to maximize the likelihood of observed text.
\end{definition}

\begin{aitip}
You can interact with a language model right now. Open a Colab notebook and run the GPT-2 example in Section~\ref{sec:first-generation}. The model will complete any text you give it.
\end{aitip}

% ============================================================
\section{Tokenization}

Language models do not process raw characters or whole words. They process \emph{tokens} --- sub-word units learned from data.

\subsection{Why sub-words?}

\begin{itemize}
  \item Character-level: vocabulary is tiny (< 300), but sequences are extremely long.
  \item Word-level: vocabulary is enormous (100\,000+), and unseen words cause failures.
  \item Sub-word: vocabulary is moderate (32\,000--100\,000), handles any word, and keeps sequences short.
\end{itemize}

\subsection{Byte Pair Encoding (BPE)}

BPE starts with individual characters and iteratively merges the most frequent adjacent pair:

\begin{enumerate}
  \item Start: \texttt{["l", "o", "w", "e", "r"]}
  \item Most frequent pair: \texttt{("l", "o")} $\rightarrow$ merge to \texttt{"lo"}
  \item Continue until vocabulary size is reached.
\end{enumerate}

\begin{minted}{python}
# Tokenization with tiktoken (GPT tokenizer)
import tiktoken

enc = tiktoken.get_encoding("cl100k_base")  # GPT-4 tokenizer
text = "Generative AI transforms how we create content."
tokens = enc.encode(text)
print(f"Text: {text}")
print(f"Token IDs: {tokens}")
print(f"Tokens: {[enc.decode([t]) for t in tokens]}")
print(f"Number of tokens: {len(tokens)}")
\end{minted}

\subsection{WordPiece tokenization}

WordPiece (used by BERT) is similar to BPE but uses a likelihood-based criterion to select merges. The HuggingFace \texttt{tokenizers} library supports both:

\begin{minted}{python}
from transformers import AutoTokenizer

# BPE tokenizer (GPT-2)
gpt2_tok = AutoTokenizer.from_pretrained("gpt2")
print(gpt2_tok.tokenize("Generative AI is transforming research."))

# WordPiece tokenizer (BERT)
bert_tok = AutoTokenizer.from_pretrained("distilbert-base-uncased")
print(bert_tok.tokenize("Generative AI is transforming research."))
\end{minted}

\begin{exercise}
\begin{enumerate}
  \item Tokenize the sentence ``Pneumonoultramicroscopicsilicovolcanoconiosis is a lung disease'' with both GPT-2 and DistilBERT tokenizers. Compare the number of tokens and the sub-word splits.
  \item Using \texttt{tiktoken}, tokenize a paragraph in English, French, and Chinese. Which language produces more tokens for roughly the same content? Why?
  \item Count the tokens in the first 1000 characters of any Wikipedia article. How does token count relate to character count?
\end{enumerate}
\end{exercise}

% ============================================================
\section{Embeddings}

Tokens are integers. Neural networks need continuous vectors. An \emph{embedding} maps each token ID to a dense vector in $\mathbb{R}^d$:

\[
\text{Embedding}: \{0, 1, \ldots, V-1\} \rightarrow \mathbb{R}^d
\]

where $V$ is the vocabulary size and $d$ is the embedding dimension (typically 768 or higher).

\begin{minted}{python}
import torch
from transformers import AutoModel, AutoTokenizer

model_name = "distilbert-base-uncased"
tokenizer = AutoTokenizer.from_pretrained(model_name)
model = AutoModel.from_pretrained(model_name)

text = "The patient has a fever."
inputs = tokenizer(text, return_tensors="pt")
with torch.no_grad():
    outputs = model(**inputs)

# outputs.last_hidden_state shape: (batch, seq_len, hidden_dim)
print(f"Shape: {outputs.last_hidden_state.shape}")
print(f"Embedding dim: {outputs.last_hidden_state.shape[-1]}")
\end{minted}

\begin{aitip}
Embeddings are not just lookup tables in modern models. After passing through Transformer layers, the output embeddings are \emph{contextual} --- the vector for ``bank'' differs depending on whether the context is finance or a river.
\end{aitip}

% ============================================================
\section{The softmax function}

The final layer of a language model produces a vector of \emph{logits} $z \in \mathbb{R}^V$ (one score per vocabulary token). The softmax converts logits to probabilities:

\[
\text{softmax}(z_i) = \frac{e^{z_i}}{\sum_{j=1}^{V} e^{z_j}}
\]

\begin{minted}{python}
import torch
import torch.nn.functional as F

# Simulated logits for a vocabulary of 5 tokens
logits = torch.tensor([2.0, 1.0, 0.5, -1.0, 3.0])
probs = F.softmax(logits, dim=0)
print("Logits:", logits.tolist())
print("Probabilities:", [f"{p:.4f}" for p in probs.tolist()])
print("Sum:", f"{probs.sum().item():.4f}")  # always 1.0
\end{minted}

% ============================================================
\section{Perplexity}
\label{sec:perplexity}

Perplexity measures how well a language model predicts a test set. Lower is better.

\[
\text{PPL} = \exp\left(-\frac{1}{N}\sum_{i=1}^{N} \log P(w_i \mid w_{<i})\right)
\]

Intuition: if perplexity is 50, the model is ``as uncertain as if it were choosing uniformly among 50 tokens'' at each step.

\begin{minted}{python}
import torch
from transformers import GPT2LMHeadModel, GPT2Tokenizer

model = GPT2LMHeadModel.from_pretrained("gpt2")
tokenizer = GPT2Tokenizer.from_pretrained("gpt2")
model.eval()

text = "The Transformer architecture revolutionized natural language processing."
inputs = tokenizer(text, return_tensors="pt")

with torch.no_grad():
    outputs = model(**inputs, labels=inputs["input_ids"])
    loss = outputs.loss  # cross-entropy loss
    perplexity = torch.exp(loss)

print(f"Loss: {loss.item():.4f}")
print(f"Perplexity: {perplexity.item():.2f}")
\end{minted}

\begin{exercise}
\begin{enumerate}
  \item Compute the perplexity of GPT-2 on three sentences: one grammatical English sentence, one scrambled version, and one sentence in French. Explain the differences.
  \item Write a function \texttt{compute\_perplexity(model, tokenizer, text)} that returns the perplexity for any input text.
\end{enumerate}
\end{exercise}

% ============================================================
\section{Your first text generation}
\label{sec:first-generation}

\begin{minted}{python}
from transformers import pipeline

generator = pipeline("text-generation", model="gpt2")
result = generator(
    "Artificial intelligence will",
    max_new_tokens=50,
    do_sample=True,
    temperature=0.8,
)
print(result[0]["generated_text"])
\end{minted}

\begin{ethicsbox}
Language models learn from internet text, which contains biases, misinformation, and harmful content. GPT-2 can generate racist, sexist, or factually wrong text. \textbf{Never deploy a base model without safety filters.} We will cover evaluation and safety in Chapter~\ref{ch:eval-safety}.
\end{ethicsbox}

% ============================================================
\section{Chapter summary}

\begin{itemize}
  \item A language model predicts the probability of the next token given previous tokens.
  \item Tokenization (BPE, WordPiece) converts text to sub-word integer sequences.
  \item Embeddings map token IDs to continuous vectors; Transformer embeddings are contextual.
  \item The softmax function converts logits to a probability distribution over the vocabulary.
  \item Perplexity measures model quality: lower perplexity means better predictions.
  \item HuggingFace \texttt{transformers} provides tokenizers, models, and generation pipelines out of the box.
\end{itemize}
